% =============================================================
% Chapter 10: The Money Behind the Machines — AI Investors and Dealmakers
% Book: AI Figures — Who's Who in the Age of AI
% =============================================================

\chapter{资本推手：AI 投资人与交易操盘者}
\label{ch:investors}

% Chapter overview: The venture capitalists, angel investors, sovereign wealth
% fund managers, and dealmakers whose capital allocation decisions are shaping
% which AI futures get built — and which do not.

如果说前面九章描绘了AI时代的技术创造者和基础设施构建者，本章要回答
一个更本质的问题：\textbf{谁在为这一切提供燃料？}

2025年，全球AI风险投资总额达到2587亿美元，占所有VC投资的61\%。
OpenAI单轮融资1100亿美元、Anthropic累计超过300亿美元、xAI一轮200亿
——这些交易的规模已超越传统风投概念，更像国家级基础设施的资本动员。
在这些天文数字背后，是一小群投资人的判断力、关系网和胆识在决定AI的
方向。他们中有些人在2019年AI尚被嘲笑为"玩具"时就下了重注，有些人则
在2024--2025年的FOMO浪潮中涌入。他们的投资逻辑、认知框架和利益纠葛，
直接塑造了我们看到的AI产业格局。

本章将全面梳理AI投资版图中最重要的人物。从硅谷的传奇风投到中国的创投
教父，从中东主权财富基金的幕后操盘手到独来独往的超级天使，从Y Combinator
的批量孵化到Stargate项目的五千亿美元豪赌——每一笔重大交易背后都有一个
具体的人做出了具体的决定。理解这些人，就是理解AI资本的运行逻辑。

\begin{whybox}[为什么AI投资人值得单独一章？]
在大多数技术叙事中，资本被简化为"背景板"——某公司融了多少钱，估值
多少。但在AI时代，投资人的角色远不止于提供资金：
\begin{enumerate}
  \item \textbf{路径选择者}：Vinod Khosla在2019年给OpenAI写下5000万美元
    支票时，OpenAI还是一个没有商业模式的非营利研究实验室。如果没有这笔
    资金，Sam Altman可能无法推动OpenAI向营利结构转型，整个GPT系列的
    发展轨迹可能完全不同。
  \item \textbf{治理参与者}：Reid Hoffman曾担任OpenAI董事，参与了多次关键
    治理决策。Josh Kushner的Thrive Capital在OpenAI的董事会安排中拥有
    重要话语权。投资人不只是签支票，他们还参与——有时甚至主导——AI公司
    的战略方向。
  \item \textbf{叙事塑造者}：Marc Andreessen的"Why AI Will Save The World"
    一文塑造了整个行业对AI风险与收益的讨论框架。a16z的Martin Casado关于
    "Bitter Economics"的分析改变了投资者对AI公司资本结构的理解。
  \item \textbf{地缘政治棋手}：MGX、GIC、QIA等主权基金的入局，让AI投资
    成为地缘政治博弈的一部分。Abu Dhabi和沙特正通过AI投资将石油财富转化
    为技术影响力。
\end{enumerate}
因此，理解这些人不仅是"八卦"，更是理解AI产业结构的必要条件。
\end{whybox}


% =============================================================
\section{十亿美元支票：谁在写AI史上最大的交易？}
\label{sec:inv-megarounds}
% =============================================================

在传统风险投资中，一笔1亿美元的交易就算"超级轮次"。但在2024--2026年的
AI融资市场中，十亿美元只是入场券。让我们从最大的交易开始，认识操盘这些
交易的关键人物。

\subsection{孙正义与SoftBank：从Vision Fund的废墟到AI的最大赌徒}
\label{subsec:inv-masa}

\begin{corebox}[人物档案：孙正义（Masayoshi Son）]
\textbf{身份}：SoftBank Group创始人兼CEO \\
\textbf{生于}：1957年，日本佐贺县鸟栖市的韩裔日本人家庭 \\
\textbf{教育}：加州大学伯克利分校经济学学士 \\
\textbf{AI关键角色}：Stargate Project联席主席；OpenAI最大外部投资者之一 \\
\textbf{投资风格}：极端集中、超高杠杆、"300年愿景"
\end{corebox}

孙正义（Masayoshi Son）是全球科技投资史上最具争议的人物之一。他的投资
生涯由两个极端定义：1990年代的雅虎和阿里巴巴投资让他一度成为世界首富，
2019--2021年Vision Fund的WeWork、Wirecard等灾难性投资则让他蒙受了
数百亿美元的亏损。但在AI时代，孙正义再次All-in。

\textbf{OpenAI豪赌。}2025年3月，SoftBank承诺向OpenAI投资高达400亿美元
——这是科技史上最大的单一投资者承诺。第一批75亿美元在2025年4月到位，
剩余225亿美元在2025年12月完成，使SoftBank在OpenAI的持股比例达到约11\%。
通过Vision Fund 2执行，SoftBank还引入了110亿美元的联合投资者。到2026年2月
OpenAI宣布1100亿美元新一轮融资时，SoftBank再次投入300亿美元——累计对
OpenAI的直接投资超过700亿美元。

\textbf{Stargate项目。}2025年1月，在白宫的一场活动中，孙正义与Sam Altman
共同宣布了"Stargate项目"——一个5000亿美元的AI基础设施联合投资计划。
参与方包括OpenAI（40\%股权）、SoftBank（40\%股权）、Oracle和MGX。
孙正义担任Stargate LLC的联席主席。SB Energy（SoftBank全资子公司）将在
德克萨斯州建设1.2吉瓦的数据中心和电力基础设施。这个项目的规模等同于
一个中等国家的能源基础设施投入。

\textbf{投资逻辑。}孙正义的AI投资逻辑可以用他自己2024年在SoftBank World
上的讲话概括："ASI（Artificial Super Intelligence）将在10年内到来，
而OpenAI是最有可能率先到达那里的公司。我要用尽一切资源，确保SoftBank
站在ASI到来时的正确位置。"这种"赌国运"式的投资风格让他在Vision Fund 1.0
时代栽了大跟头（WeWork、Oyo等），但如果OpenAI和Stargate的故事最终
成立，孙正义的AI赌注可能成为投资史上回报最高的一笔。

\begin{warnbox}[SoftBank的OpenAI投资已经赚钱了吗？]
表面上看是的。2026年2月，SoftBank在Vision Fund的季度报告中公布了42亿美元的
OpenAI投资增值。但这是纸面收益——直到OpenAI完成IPO（预计2026年底至
2027年初），SoftBank无法变现。更大的风险在于：SoftBank为了集中资源投资
AI，大幅减持了其他资产（包括部分ARM股权），杠杆率居高不下。如果AI估值
出现修正，SoftBank的整体财务状况可能迅速恶化。孙正义本人的净资产与
SoftBank股价高度绑定——这是真正的All-in。
\end{warnbox}


\subsection{Josh Kushner与Thrive Capital：OpenAI最亲密的金融伙伴}
\label{subsec:inv-thrive}

\begin{corebox}[人物档案：Josh Kushner]
\textbf{身份}：Thrive Capital创始人兼管理合伙人 \\
\textbf{生于}：1985年，美国新泽西州 \\
\textbf{教育}：哈佛大学 \\
\textbf{AI关键角色}：OpenAI最早期的机构投资者之一；多轮领投或重要参投 \\
\textbf{投资风格}：成长期、收入驱动、长期伙伴关系 \\
\textbf{净资产}：估计超过40亿美元
\end{corebox}

Josh Kushner在2011年创立Thrive Capital时只有26岁。过去十五年中，Thrive
投资了Instagram、Spotify、Stripe等定义时代的科技公司。但真正让Kushner
跻身AI投资最核心圈层的是他与Sam Altman的关系，以及对OpenAI的持续加注。

\textbf{OpenAI投资时间线：}
\begin{itemize}
  \item \textbf{2023年}：Thrive领导了一轮OpenAI二级市场交易，估值860亿美元
    ——这是Kushner首次大规模介入OpenAI；
  \item \textbf{2024年10月}：Thrive参与了OpenAI 66亿美元融资轮，
    投资约13亿美元，Kushner公开表示"因为两个关键原因下注"；
  \item \textbf{2025年3月}：参与OpenAI 400亿美元轮次；
  \item \textbf{2025年12月}：以2850亿美元估值单独投资约10亿美元
    ——当时OpenAI已在讨论以8300亿美元估值的更大轮次，这意味着
    Thrive的这笔投资一入账就"已经在赚钱"（in the money）；
  \item \textbf{2026年2月}：预计参与1100亿美元轮次。
\end{itemize}

Kushner的成功不仅在于投资眼光，更在于关系经营。Sam Altman对Kushner的
评价是："Extremely grateful to work with Josh. No one could ask for a more
committed, more thoughtful, or harder-working investor."这种密切关系让
Thrive获得了其他基金难以比拟的优先配额权——在AI融资市场中，
\textbf{拿到配额往往比选对公司更难}。

2026年2月，Thrive宣布关闭了其最新基金，规模超过100亿美元——这使
Thrive成为仅有的几家单只基金突破百亿的成长期基金之一。OpenAI也
反向入股了Thrive，双方的利益深度绑定。

\begin{whybox}[为什么Thrive能获得OpenAI的"特殊交易"？]
2025年12月，Thrive以2850亿美元估值投资OpenAI——几个月后OpenAI就以
8400亿美元估值融资。这种"优惠价"引发了市场争议。《华尔街日报》报道，
这种不同投资者获得不同条件的做法在AI融资中越来越常见——核心原因是
AI公司需要的资金规模如此之大，以至于它们需要"锚定投资者"来建立市场
信心。Thrive作为OpenAI最长期的财务投资者，扮演了这个角色。代价是
什么？后进投资者可能需要接受更高的估值和更弱的保护条款——这实际上是
用"不平等"的交易条件来奖励忠诚度和长期承诺。
\end{whybox}


% =============================================================
\section{硅谷AI权力掮客}
\label{sec:inv-sv-power}
% =============================================================

如果说孙正义和Kushner代表了"大支票"投资者，本节介绍的人物则代表了
硅谷AI投资的思想领导力——他们不仅投资金，更投信念和愿景。

\subsection{Marc Andreessen：AI乐观主义的旗手}
\label{subsec:inv-marc}

\begin{corebox}[人物档案：Marc Andreessen]
\textbf{身份}：Andreessen Horowitz（a16z）联合创始人兼普通合伙人 \\
\textbf{生于}：1971年，美国爱荷华州 \\
\textbf{教育}：伊利诺伊大学厄巴纳-香槟分校计算机科学学士 \\
\textbf{早期成就}：Mosaic浏览器创建者；Netscape联合创始人 \\
\textbf{AI关键角色}：a16z AI投资战略总设计师；AI政策倡导者 \\
\textbf{代表性文章}："Why AI Will Save The World"（2023年6月）
\end{corebox}

Marc Andreessen可能是当代科技投资界最具思想影响力的人物。他在1993年
共同创建了Mosaic浏览器（后来的Netscape），直接催生了互联网时代。
2009年与Ben Horowitz共同创立a16z后，他提出的"Software Is Eating The
World"（2011年《华尔街日报》专栏）成为过去十年科技投资的核心论述。
进入AI时代，Andreessen再次扮演了叙事塑造者的角色。

\textbf{``Why AI Will Save The World''——AI时代的宣言书。}2023年6月，
Andreessen在a16z网站上发表了长文"Why AI Will Save The World"，系统性地
反驳了AI威胁论，主张AI将"增强而非取代人类智能"，并将AI监管描述为
"既有利益集团保护护城河的工具"。这篇文章迅速成为AI投资圈的"必读文本"
——无论你是否同意其观点，它为AI乐观主义提供了一个完整的论证框架。

\textbf{a16z的AI投资版图。}在Andreessen的战略指导下，a16z在AI领域的
部署堪称全方位：
\begin{itemize}
  \item OpenAI——据报道投资超过100亿美元；
  \item Anthropic——据报道投资约40亿美元；
  \item Databricks, Scale AI——数据基础设施层；
  \item Cursor（代码AI，估值500亿美元的开发者工具）；
  \item ElevenLabs——语音合成AI领头羊；
  \item Character.AI——消费者AI应用；
  \item Mistral AI——欧洲最重要的基础模型公司。
\end{itemize}

a16z的管理资产（AUM）已超过900亿美元，使其成为全球最大的科技风投基金。
2025年，a16z在AI领域参与了超过25笔交易，单笔投资规模从500万到5亿美元
不等。

2026年1月，Andreessen在播客中进一步阐述了他对AI时代的看法："2025年
可能是科技史上最重要的一年……这不是渐进式改进，而是从零到一的转变。
AI正在重新定义产品经理、设计师和工程师的角色。"

\textbf{政治参与。}Andreessen在2024年美国大选中公开支持特朗普，
这一举动在硅谷引发了强烈争议。他的政治立场与AI政策主张紧密关联：
反对过度监管、支持"Little Tech"、反对FARA等对创业公司构成合规负担
的法规。他与Ben Horowitz共同推出的"The Little Tech Agenda"将AI监管
批评扩展为更广泛的科技政策框架。

\begin{keybox}[a16z AI团队核心人物]
\small
\begin{tabularx}{\textwidth}{l l X}
\toprule
\rowcolor{figLightGray}
\textbf{姓名} & \textbf{角色} & \textbf{关键贡献} \\
\midrule
Marc Andreessen & 联合创始人 & AI战略总设计师，"Why AI Will Save The World" \\
Ben Horowitz    & 联合创始人 & a16z文化构建者，2024年政治参与 \\
Martin Casado   & 普通合伙人 & AI基础设施投资论文，"Bitter Economics"分析 \\
Anjney Midha    & 风险合伙人 & GPU分配决策者，Mistral/BFL董事，AMP创始人 \\
Matt Bornstein  & 合伙人 & AI应用层投资，AI基础设施堆栈分析 \\
\bottomrule
\end{tabularx}
\end{keybox}

\paragraph{Martin Casado——AI基础设施的思想家}

Martin Casado在2015年以普通合伙人身份加入a16z之前，是SDN（软件定义
网络）领域的先驱——他创立的Nicira在2012年被VMware以12.6亿美元收购。
在a16z，Casado成为AI基础设施投资的核心驱动力。

2025年11月，Casado发表了极具影响力的文章"Bitter Economics"，将Richard
Sutton的"The Bitter Lesson"（计算胜过巧妙算法）扩展到商业层面：
AI系统本质上是"将大量资本转化为可工作的解决方案"，这从根本上改变了
软件公司的资本结构。传统软件公司的边际成本趋近于零，但AI公司的边际
成本（推理算力）始终存在——这意味着AI公司更像资本密集型企业而非
软件公司。这一洞见深刻影响了投资者对AI公司估值的理解。

2026年1月，Casado在播客中进一步解释了AI需求侧的驱动力：当前的约束
更多来自监管（电力、数据中心、算力扩张方面的限制），而非技术本身。
AI Agent工具可能将基础设施决策从人类转移到智能体——这将是基础设施
投资的下一个重大转折点。

\paragraph{Anjney Midha——算力的守门人}

Anjney Midha的职业轨迹揭示了AI时代一个鲜为人知但极其重要的权力节点：
\textbf{算力分配权}。

Midha在2023年加入a16z后，负责管理公司的"Oxygen"项目——一个超过
20000块GPU的私有算力集群。在GPU严重短缺的2023--2024年，Midha决定哪些
AI创业公司能获得a16z的算力支持——这实质上是在决定哪些公司能生存。
他同时担任Mistral AI（欧洲最大的开源LLM公司，估值140亿美元）和Black
Forest Labs（德国图像生成AI公司，估值32.5亿美元）的董事。

2024年10月，Midha宣布创立AMP——一家独立的"算力+资本"公司，为前沿AI
团队同时提供GPU资源和投资。他继续以风险合伙人身份留在a16z。这一举动
暴露了AI时代的一个根本真相：\textbf{在算力成为稀缺资源的时代，算力
分配者拥有的权力可能不亚于资本分配者}。Midha年仅33岁，已经坐在了
AI世界的一个关键交叉点上。


\subsection{Vinod Khosla：AI史上最成功的单笔投资}
\label{subsec:inv-khosla}

\begin{corebox}[人物档案：Vinod Khosla]
\textbf{身份}：Khosla Ventures创始人 \\
\textbf{生于}：1955年，印度新德里 \\
\textbf{教育}：IIT Delhi电气工程学士；卡内基梅隆大学生物医学工程硕士；
  斯坦福大学商学院MBA \\
\textbf{早期成就}：Sun Microsystems联合创始人 \\
\textbf{AI关键投资}：OpenAI（2019年5000万美元→2026年价值约80亿美元，160倍回报） \\
\textbf{投资风格}："黑天鹅猎人"——押注科学实验级别的颠覆性技术
\end{corebox}

Vinod Khosla对OpenAI的投资可能是风险投资史上最成功的单笔交易之一。

\textbf{故事从Elon Musk的退出开始。}2019年，OpenAI正处于生死关头。
作为非营利组织成立四年后，OpenAI面临一个根本性矛盾：深度学习研究需要
海量算力，但非营利架构限制了其筹资能力。联合创始人Elon Musk在2018年
退出董事会后，还留下了对10亿美元资金承诺的不确定性。Sam Altman需要
找到新的资金来源来支撑OpenAI的野心。

Khosla此时入场。据Fortune报道，在Musk犹豫之际，Khosla果断写下了5000万
美元的支票——这是他40年投资生涯中最大的单笔初始投资。The Information
报道，这笔投资换取了OpenAI约5\%的股权。到2026年初，以OpenAI 8400亿美元
的估值计算，这笔投资价值约80亿美元——回报率约160倍。

\textbf{这不是运气，是认知框架。}Khosla的投资哲学可以用"指数套利"
（Exponential Arbitrage）来概括：当技术进步的速度是指数级的，而大多数人
的认知仍然是线性的，指数与线性之间的差距就是巨大的投资机会。2019年的
OpenAI正处于这样一个认知差距中——GPT-2已经展示了语言模型的惊人潜力，
但绝大多数投资者仍然认为AI是"学术玩具"。

Khosla在2019年看到的，是GPT-2到GPT-3的"跃迁"——同样的架构，更多的
算力和数据，能力呈指数级增长。他赌的不是某个具体产品，而是scaling laws
——算力投入和模型能力之间的幂律关系会持续成立。事实证明，他对了。

\textbf{持续加注。}Khosla不仅在2019年下了初始注，还在后续轮次中持续
加注。2024年10月，Khosla Ventures通过SPV（特殊目的载体）为OpenAI的66亿
美元轮次额外筹集了4.05亿美元。值得注意的是，这4.05亿中的大部分来自
外部投资者——Khosla利用自己的配额权帮助其他渴望参与OpenAI的投资者
获得入场券。这种模式在AI融资中越来越常见：顶级VC的核心价值不仅在于
自己的资金，更在于\textbf{分配稀缺配额的权力}。

\textbf{AI预言。}Khosla是AI领域最大胆的预言者之一。他公开预测到2030年
80\%的现有工作岗位将被AI取代——同时强调这将创造"前所未有的繁荣"。
2026年3月，他在Fortune的采访中讨论了谁在AI竞赛中领先，表达了对OpenAI
保持领先地位的信心，同时也强调了美国需要加速AI发展以在全球竞争中
保持优势。


\subsection{Reid Hoffman：AI生态系统的连接者}
\label{subsec:inv-hoffman}

\begin{corebox}[人物档案：Reid Hoffman]
\textbf{身份}：Greylock Partners合伙人；LinkedIn联合创始人 \\
\textbf{生于}：1967年，美国加利福尼亚州 \\
\textbf{教育}：斯坦福大学符号系统学士；牛津大学硕士 \\
\textbf{AI关键角色}：OpenAI早期董事会成员和投资者；Anthropic投资者；
  AI政策倡导者 \\
\textbf{标签}：``PayPal Mafia''成员；``硅谷的外交官''
\end{corebox}

Reid Hoffman的独特之处在于，他可能是AI领域中同时押注了最多"竞争对手"
的投资者。作为个人投资者和Greylock合伙人，Hoffman同时投资了OpenAI和
Anthropic——两家正面竞争的AI公司。

\textbf{OpenAI的早期支持者。}Hoffman是OpenAI最初的捐赠者和董事会成员之一，
参与了OpenAI从非营利到营利结构转型的关键决策。他在OpenAI董事会一直任职
到2023年，期间经历了包括Sam Altman短暂被解雇又复职的戏剧性事件。

\textbf{同时投资Anthropic。}2025年10月，当AI和加密沙皇David Sacks公开
批评Anthropic时，Hoffman公开为Anthropic辩护，称其为"好人之一"（one of
the good guys），并表示："Anthropic与Microsoft、Google和OpenAI一样，正在
以正确的方式部署AI——深思熟虑、安全且对社会有巨大益处。"他同时确认
Greylock已投资Anthropic。

这种"两边下注"的策略在传统VC中会引发利益冲突质疑，但在AI领域却越来越
普遍——因为投资者的共识是：AI市场足够大，容得下多个赢家。Hoffman的
核心信念是AI本身的转型力量，而非某家特定公司的胜出。

\textbf{思想影响力。}Hoffman是AI领域最多产的公共知识分子之一。他共同
撰写了多本书，包括与GPT-4合作撰写的《Impromptu: Amplifying Our Humanity
Through AI》。他的播客和专栏持续探讨AI对社会、就业和民主的影响。
与Andreessen的纯乐观主义不同，Hoffman更倾向于"cautious optimism"
——承认AI的风险，但相信通过负责任的开发可以最大化收益。


\subsection{Sequoia Capital：AI投资的系统化玩家}
\label{subsec:inv-sequoia}

\begin{corebox}[Sequoia Capital AI投资概览]
\textbf{管理资产}：超过900亿美元 \\
\textbf{2025年AI交易数}：20+ \\
\textbf{代表性AI投资}：Anthropic, xAI, Scale AI, Glean, ElevenLabs,
  Hugging Face, Mistral AI \\
\textbf{2026年AI VC排名}：顶级（Earthian AI、PitchBook等多个排名）
\end{corebox}

\paragraph{Roelof Botha——稳健的舵手}

Roelof Botha自2003年加入Sequoia以来，逐步成为硅谷最受尊敬的投资人之一。
出生于南非，Botha早年在PayPal担任CFO（与Elon Musk、Peter Thiel共事），
之后加入Sequoia。他领导了对YouTube、Instagram、MongoDB等公司的早期投资。

在AI时代，Botha带领Sequoia迅速布局——Anthropic的多轮融资、xAI、Scale AI
等关键交易都有Sequoia的身影。2025年11月，Botha在任职近十年后将Sequoia
的"管家"（steward）角色移交给Pat Grady和Alfred Lin。Bloomberg报道，
这一决定是合伙人一致做出的，Grady和Lin计划在Botha的基础上"深化对AI的
聚焦"，同时降低公司的政治化形象。

\paragraph{Pat Grady——Sequoia的AI操盘手}

Pat Grady是Sequoia AI战略的实际操盘者。2007年，24岁的Grady加入Sequoia
后，靠的不是天才的投资直觉，而是"笨功夫"——每天打50个陌生电话
（cold calls），系统性地扫描潜在投资标的。他在2024年公开分享了Sequoia
早年投资ServiceNow的CRM记录："left message with assistant"、
"sent email"、"left voice mail"——一周又一周的坚持，直到Fred Luddy
终于回电。

Grady与另一位合伙人Sonya Huang共同在Sequoia官网上发表公开信，邀请
AI创始人直接向他们发送想法和pitch——这种"开放式招商"的方式在顶级VC中
极为罕见。Business Insider将Grady列入2024年AI Power List，评价他
"比任何人都更深入地领导了Sequoia的AI战略"。

2025年11月，43岁的Grady成为Sequoia的联席管家（co-steward）。他的升任
标志着Sequoia明确将AI作为未来十年的核心投资方向。

\begin{whybox}[Sequoia对AI交易的态度：不追求好价格]
Grady曾对Business Insider说："We don't make money by getting a good deal."
——Sequoia的成功不是靠在估值上占便宜，而是靠投对人。这种哲学在AI时代
尤为重要：当OpenAI估值8400亿、Anthropic估值3800亿时，纠结于某轮融资的
估值差异几乎没有意义——真正的问题是这些公司能否在未来十年创造数千亿
美元的收入。如果答案是肯定的，那么在1000亿还是1200亿估值时入场的差异
微不足道。
\end{whybox}


\subsection{超级天使与双重身份投资者}
\label{subsec:inv-angels}

AI时代产生了一类独特的投资者：他们既是技术从业者（创始人、工程师、
高管），又是活跃的天使投资人。他们的投资判断往往源于对技术的第一手理解，
而非传统的财务分析。

\paragraph{Daniel Gross与Nat Friedman：从投资人到被收购者}

Daniel Gross和Nat Friedman的故事是AI时代最戏剧性的投资人叙事之一。

Gross早年在Apple担任AI和机器学习主管，之后创立了Pioneer（一个在线
创业加速器）。Friedman则是GitHub的前CEO，在Microsoft收购GitHub的过程中
发挥了关键作用。两人于2021年联合创立了NFDG（Nat Friedman \& Daniel Gross），
一个专注AI的风投基金，投资了Perplexity、The Bot Company等热门公司。

2024年6月，Gross成为Safe Superintelligence（SSI）的联合创始人兼CEO——
SSI由OpenAI联合创始人Ilya Sutskever发起，专注于开发"安全的超级智能"，
估值迅速攀升至320亿美元。这使Gross从投资者变成了创始人。

然后事情发生了戏剧性转折。2025年6月，Meta试图以320亿美元收购SSI
——被Sutskever拒绝。收购失败后，Zuckerberg改为直接挖人：Meta聘请了
Gross和Friedman，同时部分收购了NFDG基金的份额，获得了该基金所投AI
创业公司的少数股权。这一系列操作展示了AI时代的一个新现象：
\textbf{投资人本身成为了被收购的目标}——因为他们的关系网、技术洞察力
和所持有的投资组合，比任何单笔投资都更有价值。

\paragraph{Elad Gil：一个人的多十亿帝国}

Elad Gil可能是AI时代最成功的solo investor（单人投资者）。他的职业轨迹
从Google早期员工（负责Google Maps和Gmail等移动产品）到Twitter高管
（通过收购其创立的Mixer Labs加入），再到全职天使投资人。

2024年，Gil筹集了10亿美元的新基金，使其个人可投资资金超过20亿美元。
他的AI投资组合涵盖了几乎所有热门赛道：
\begin{itemize}
  \item \textbf{基础模型}：OpenAI, Mistral；
  \item \textbf{AI应用}：Harvey（法律AI）, Pika（视频生成）,
    Perplexity（AI搜索）；
  \item \textbf{AI基础设施}：Magic（代码生成）；
  \item \textbf{防务AI}：Helsing（欧洲AI防务独角兽）。
\end{itemize}

2025年11月在TechCrunch Disrupt大会上，Gil分享了他的AI投资哲学：
"AI是我见过的最不可预测的技术浪潮之一……我在2021年开始投资生成式AI，
当时几乎没人关注。GPT-2到GPT-3的跃迁如此巨大，如果你只是沿着
scaling laws外推，就能看到这将变得极其重要。"

Gil的投资风格有一个显著特征：\textbf{当他出现在cap table上时，估值
立刻上升}。更大的投资者知道，如果Elad已经投了，那这家公司值得认真
看——这是一种强大的信号效应。

\paragraph{Naval Ravikant：哲学投资家}

Naval Ravikant作为AngelList的创始人，不仅是AI投资者，更是整个天使投资
基础设施的构建者。AngelList使得小额天使投资人能够高效地参与创业公司
融资——在AI时代，数百个AI创业公司通过AngelList的Rolling Funds和
Syndicates获得了种子资金。

Ravikant的投资哲学更接近于一种技术哲学：他在Twitter上广泛讨论AI对
人类认知、工作和意义的影响。他的核心观点是AI将使"杠杆"（leverage）
民主化——以前只有资本和劳动力能产生杠杆，现在代码和AI为每个人提供了
无限杠杆。这种思想框架影响了一整代AI创业者的自我认知和创业方向。

\paragraph{Eric Schmidt：从Google CEO到AI战略家}

Eric Schmidt的身份转变可能是科技史上最引人注目的之一——从全球最大
AI公司之一（Google）的掌门人，到独立的AI投资者、政策倡导者和创业者。

\textbf{投资与创业。}Schmidt离开Google后，通过个人投资和Schmidt Futures
（其慈善科技基金）广泛投资AI领域。他还共同创立了Relativity Space，
并担任CEO。但他最大的影响力来自\textbf{政策层面}：他曾领导美国国家
人工智能安全委员会（NSCAI），该委员会在2021年发布的报告深刻影响了
美国的AI战略方向。

\textbf{AI地缘政治论述。}2025年12月，Schmidt在Time杂志发表文章
"How the U.S. Can Win the AI Race"，提出AI竞赛实际上是多个重叠的竞赛
——美中AGI竞争、AI应用普及竞争、AI治理竞争。他指出中国的AI+计划
目标到2030年实现90\%的关键行业AI整合，而美国在AI采用率上反而落后于
中国（中国60\%的员工每周使用AI，美国仅37\%）。

2026年3月，Schmidt在Fortune发表专栏，主张科技巨头应该建设与电力源共址
的数据中心（co-located data centers），而非依赖公共电网——直接回应了
公众对AI数据中心推高电费的担忧。这种从投资者到政策倡导者的角色转换，
使Schmidt成为AI时代最独特的"混合身份"人物之一。

\paragraph{Jeff Bezos：低调的AI多头}

Jeff Bezos可能是AI领域投资规模最大但最低调的个人投资者。通过Bezos
Expeditions（其个人投资工具），Bezos在OpenAI成立初期就参与了资助
——他是2015年最早向OpenAI捐款的科技领袖之一。

但Bezos更大的AI投资是通过Amazon完成的——Amazon对Anthropic的累计投资
承诺达到80亿美元。到2025年2月，Amazon在Anthropic的持股价值已升至约
140亿美元，账面增值约60亿美元（75\%回报率）。此外，Amazon在2026年2月
以500亿美元参与了OpenAI的1100亿美元融资——标志着Amazon从"Anthropic独家
盟友"转变为"AI双头下注"的策略。

\paragraph{Peter Thiel与Founders Fund：从谨慎到All-in}

Peter Thiel的AI投资之路经历了一个有趣的转变。Founders Fund最初对
生成式AI持"谨慎观望"态度——这对一家以"Contrarian"著称的基金来说
并不意外。但到2025年下半年，Founders Fund转向了"集中押注"模式。

Founders Fund的AI投资组合反映了Thiel一贯的"与国家竞争"
（competing with nation-states）论述：
\begin{itemize}
  \item \textbf{Anthropic}：参与了\$30B Series G融资的联合领投方之一；
  \item \textbf{Anduril}：AI防务系统公司，\$25亿Series G；
  \item \textbf{Scale AI}：AI数据基础设施，\$10亿Series F；
  \item \textbf{Crusoe Energy}：AI数据中心能源公司，\$5亿融资；
  \item \textbf{Hadrian}：AI驱动的精密制造。
\end{itemize}

2026年3月，Founders Fund正在筹集60亿美元的新基金（Growth IV），预计
将大幅加重AI在投资组合中的比重。Thiel的AI投资论述可以概括为：
\textbf{AI不是一般意义上的技术投资，而是关乎美国国家安全的战略资产}。
这种框架与他在Palantir的创始经历和长期的国防投资偏好一脉相承。


\subsection{Y Combinator：AI创业的批量生产线}
\label{subsec:inv-yc}

\begin{corebox}[人物档案：Garry Tan]
\textbf{身份}：Y Combinator CEO \\
\textbf{生于}：1981年，加拿大温尼伯，华裔 \\
\textbf{教育}：斯坦福大学计算机系统工程学士 \\
\textbf{早期经历}：Palantir第10号员工；Posterous联合创始人
  （被Twitter收购）；Initialized Capital联合创始人 \\
\textbf{AI关键角色}：重塑YC为AI创业首选加速器
\end{corebox}

Garry Tan于2023年接任Y Combinator CEO后，迅速将这个传统的创业加速器
转变为AI创业的核心枢纽。在他的领导下，YC每个批次中AI公司的比例从2023年
的约32\%飙升至2025年的67\%以上。

\textbf{AI改变了YC的选人标准。}2026年3月在SXSW的演讲中，Tan透露AI正在
根本性地改变YC评估申请者的方式。当编程智慧"随手可得"（on tap），加速器
更看重创始人的\textbf{品味、行动力和产品管理能力}。"下一个伟大的发明者
可能是英语专业的学生"——这句话标志着创业门槛的范式转变。

Tan本人也是AI工具的狂热用户。他在SXSW上对Bill Gurley说自己只睡四个
小时，因为太兴奋于和AI Agent协作。"我可以做我的全职工作——八九个小时
的会议——同时在三个不同项目上完成10000行代码。"他开玩笑说自己得了
"cyber psychosis"（AI引起的兴奋失眠）。

YC的每个批次为每家公司提供50万美元的种子资金和13周的加速指导。更重要的是
YC的网络效应：Airbnb、DoorDash、Coinbase、Instacart等超级独角兽都出自
YC——这种校友网络为AI创业公司提供了客户、合作伙伴和后续融资的巨大优势。


% =============================================================
\section{中国AI投资版图}
\label{sec:inv-china}
% =============================================================

中国的AI投资生态与硅谷有着根本性的差异：更深的政府参与、更紧密的产业-
资本联动、以及地缘政治对跨境投资的严重制约。在这个生态中，一小群投资人
扮演着特殊的"翻译者"角色——在中国和全球科技之间架桥。

\subsection{沈南鹏（Neil Shen）：两个世界之间的桥梁}
\label{subsec:inv-shen}

\begin{corebox}[人物档案：沈南鹏（Neil Shen）]
\textbf{身份}：红杉中国/HongShan（HSG）创始及管理合伙人 \\
\textbf{生于}：1967年，浙江海宁 \\
\textbf{教育}：上海交通大学应用数学学士；耶鲁管理学院硕士 \\
\textbf{早期成就}：携程旅行网（Ctrip）联合创始人；如家酒店连锁联合创始人 \\
\textbf{AI关键角色}：中国AI投资规模最大、影响力最深的单一投资人 \\
\textbf{净资产}：约38亿美元（2025年9月Forbes估算） \\
\textbf{管理资产}：约560亿美元
\end{corebox}

沈南鹏是中国风险投资史上最成功的人物——这一点得到了Forbes和
Financial Times的一致认可。他在2005年创立红杉资本中国后，投资了
包括美团、字节跳动、拼多多等定义中国互联网时代的公司。

\textbf{从红杉中国到HongShan的独立。}2023年，全球红杉拆分为三个独立
实体——美国Sequoia Capital、中国HongShan（原Sequoia China）和
印度/东南亚Peak XV。沈南鹏领导的HongShan成为完全独立的投资机构。这次
拆分反映了中美地缘政治紧张的深远影响：美国LP（有限合伙人）越来越不愿
让自己的资金流入中国科技公司，而中国政府也在审视美资对国内科技企业的
影响力。

\textbf{AI投资布局。}在AI时代，沈南鹏的HongShan依然是中国AI领域最重要
的资本力量。2026年2月，《华尔街日报》以"将美国资本注入中国AI的香港
投资人"（The Hong Kong Investor Putting American Money Into China's AI
Push）为题报道了沈南鹏如何继续利用美国资本支持中国AI公司——尽管
面临越来越大的地缘政治压力。

HongShan在AI领域的投资覆盖了中国AI生态的几乎所有关键层：
\begin{itemize}
  \item \textbf{基础模型}：月之暗面（Moonshot AI）、智谱AI（Zhipu AI）等
    中国头部大模型公司；
  \item \textbf{AI应用}：多个垂直领域的AI应用创业公司；
  \item \textbf{AI基础设施}：芯片、算力、数据相关的创业公司。
\end{itemize}

沈南鹏的核心挑战在于：如何在中美技术脱钩的趋势下，继续维持中国AI
创业公司获取全球资本和技术的通道？HongShan将总部设在香港而非北京或
上海，本身就是一种策略性安排——香港作为中间地带，为跨境资本流动提供
了一个虽然收窄但仍然存在的通道。


\subsection{张磊（Zhang Lei）：从耶鲁到中国AI工厂}
\label{subsec:inv-zhanglei}

\begin{corebox}[人物档案：张磊（Zhang Lei）]
\textbf{身份}：高瓴投资（Hillhouse Investment）创始人兼CEO \\
\textbf{生于}：1972年，中国河南省驻马店市 \\
\textbf{教育}：中国人民大学学士；耶鲁管理学院MBA \\
\textbf{早期成就}：2005年用耶鲁大学捐赠基金提供的2000万美元创立高瓴 \\
\textbf{管理资产}：超过150亿美元（公开市场合并基金） \\
\textbf{AI观点}："中国已成为世界的AI工厂"
\end{corebox}

张磊和高瓴在中国投资界的地位类似于Sequoia在硅谷的角色——以长期主义和
研究驱动著称。高瓴早期以投资腾讯和京东而闻名，管理过超过180亿美元的
旗舰私募基金。

\textbf{AI战略定位。}在AI时代，张磊的策略与多数纯财务投资者不同——
他更强调"AI+产业"的融合。2025年10月，张磊在Global Financial Leaders'
Investment Summit上表示："中国将最先在AI应用层实现大规模变现。"
他的论点是：中国虽然在基础模型的绝对能力上可能暂时落后于美国
（OpenAI、Anthropic、Google的模型在某些benchmark上领先），但在
\textbf{AI应用的速度和广度}上，中国已经领先。

2025年9月，高瓴的风险投资部门GL Ventures与浦东创投共同成立了20亿元人民币
（约2.8亿美元）的AI专项基金，支持上海张江AI创新小镇的建设，目标是
孵化1000家AI公司。这种"产业+资本+地方政府"的三角合作模式是中国AI投资
的典型特征——与硅谷纯市场驱动的模式形成鲜明对比。

2025年11月，张磊将高瓴旗下三只公开市场基金合并为一只超过150亿美元的
综合基金——这是多年来最大规模的内部重组，旨在让投资组合更加灵活。
在AI领域，高瓴的投资覆盖了自动驾驶（地平线机器人）、AI制造、AI医疗等
"重型应用"赛道。


\subsection{李开复（Kai-Fu Lee）：科学家、投资人、创业者的三重身份}
\label{subsec:inv-kaifulee}

\begin{corebox}[人物档案：李开复（Kai-Fu Lee）]
\textbf{身份}：创新工场（Sinovation Ventures）董事长兼CEO；01.AI创始人兼CEO \\
\textbf{生于}：1961年，台湾台北 \\
\textbf{教育}：哥伦比亚大学计算机科学学士；卡内基梅隆大学计算机科学博士 \\
\textbf{职业路径}：Apple VP → SGI VP → Microsoft Corp VP →
  Google中国区总裁 → 创新工场创始人 → 01.AI创始人 \\
\textbf{代表作}：《AI Superpowers》《AI 2041》 \\
\textbf{管理资产}：约30亿美元（双币种基金）
\end{corebox}

李开复是AI领域罕见的"三栖人物"——他同时是顶级的AI研究者（博士论文创建了
世界上第一个不依赖特定说话人的连续语音识别系统SPHINX）、活跃的投资者
（创新工场管理30亿美元基金）和一线创业者（01.AI）。

\textbf{从投资者到创业者。}2023年，李开复做出了令投资圈惊讶的决定：
在62岁时亲自创立01.AI，一家专注于领域特定大语言模型的生成式AI公司。
01.AI迅速成为独角兽，其开源模型Yi系列在特定benchmark上表现出色。
TIME将他列入"Top 25 AI Leaders"，世界经济论坛邀请他担任AI Council
联席主席。

\textbf{中美AI差距的权威解读者。}李开复可能是最有资格同时理解中美AI
生态的人——他在Apple、Microsoft和Google（均在中国有深度运营）担任过
高级管理职位，同时又是中国本土AI创业生态的深度参与者。

2025年11月，他在Bloomberg的采访中表示："美国公司将率先实现AI服务的
商业化变现，之后中国公司将紧随其后。"这种"时间差"论述为中国AI投资者
提供了一个框架：不需要在基础模型上与OpenAI正面竞争，而是在
\textbf{应用层和本地化}上寻找独特的赢面。

\textbf{创新工场的AI布局。}作为投资者，李开复通过创新工场投资了大量
中国AI创业公司。创新工场的特色在于其"AI研究院"——一个内部的技术研发
团队，不仅投资AI公司，还直接进行AI研究和人才培养。这种"VC+研究院"
的混合模式在全球风投中几乎独一无二。


\subsection{朱啸虎（Allen Zhu）：实用主义的AI投资者}
\label{subsec:inv-zhu}

\begin{corebox}[人物档案：朱啸虎（Allen Zhu Xiaohu）]
\textbf{身份}：金沙江创投（GSR Ventures）管理合伙人 \\
\textbf{教育}：上海交通大学；麦肯锡大中华区工作经验 \\
\textbf{早期成就}：美团（Meituan）和滴滴出行（Didi Chuxing）的早期投资者 \\
\textbf{投资风格}："狙击手型"——少量投资、深度研究、强执行导向 \\
\textbf{AI观点}：强调用户留存和商业化能力，对人形机器人持怀疑态度
\end{corebox}

朱啸虎是中国VC界最直言不讳的声音之一。在AI投资热潮中，他的观点常常
与主流"FOMO"情绪形成对比。

\textbf{AI投资哲学。}2025年9月在外滩大会上，朱啸虎提出了一个简洁但
深刻的AI投资框架："用户留存率是投资者评估AI创业公司时看的最重要的
单一指标。"他指出AI时代的一个普遍问题：用户喜欢尝试新产品，但第一个
月付费后第二个月就流失。"一旦用户流失，拉回来的成本是获客成本的10倍
以上，几乎不可能。"

\textbf{对人形机器人的公开质疑。}2025年3月，朱啸虎在《中国企业家》的
采访中公开质疑了人形机器人的商业可行性——这在当时中国Embodied AI投资
热潮的背景下显得格格不入。他透露GSR Ventures已经退出了多个Embodied AI
投资，理由是"商业化路径不清晰"。这一观点在中国机器人圈引发了热烈辩论。

朱啸虎的投资风格被业界称为"sniper investor"（狙击手投资者）——他不追求
投资数量，而是对每个项目进行深入研究后精准下注。在AI时代，他更倾向于
投资有明确收入模式和用户粘性的\textbf{应用层}公司，而非烧钱的基础模型。


\subsection{陆奇（Lu Qi）：中国AI创业的教父}
\label{subsec:inv-luqi}

\begin{corebox}[人物档案：陆奇（Lu Qi）]
\textbf{身份}：奇绩创坛（MiraclePlus）创始人 \\
\textbf{生于}：1961年，中国上海 \\
\textbf{教育}：复旦大学计算机科学学士和硕士；卡内基梅隆大学计算机科学博士 \\
\textbf{职业路径}：Yahoo EVP → Microsoft EVP（直接向Satya Nadella汇报）
  → 百度COO兼总裁 → YC中国创始人 → 奇绩创坛创始人 \\
\textbf{AI关键角色}：中国版Y Combinator的缔造者
\end{corebox}

陆奇的职业生涯本身就是一部中美AI交叉史。他在Microsoft时期是Satya
Nadella最信任的副手之一，负责Applications and Services Group（包括Office、
Bing等）。2017年，他以百度COO身份回到中国，被视为百度AI转型的关键推手
——但仅14个月后就因"个人原因"离职。

\textbf{从YC中国到奇绩创坛。}2018年，陆奇受邀创立Y Combinator中国分部。
但2019年受中美贸易战影响，YC宣布关闭中国分部。陆奇随即将团队独立，
创立了奇绩创坛（MiraclePlus）——本质上是中国版的YC。

奇绩创坛每年举办两次创业路演（春季和秋季），主要面向IT和深科技领域的
早期创业者。2024年，奇绩创坛进行了12笔公开投资，覆盖了77家企业。投资
方向高度集中在AI相关领域：机器人、自动驾驶、空间智能、AI应用等。
其投资的创业者典型画像是：名校毕业、25--30岁、在顶级公司有几年工作
经验的年轻技术人才。

陆奇在奇绩创坛中扮演的角色不仅仅是投资者——他更像是中国AI创业者的
\textbf{导师和精神领袖}。他在每次路演中的点评和指导被创业者广泛引用，
其对AI技术趋势的判断在中国创投圈具有极大影响力。


% =============================================================
\section{主权资本与AI地缘政治}
\label{sec:inv-sovereign}
% =============================================================

AI投资的第三极力量不是来自硅谷的风投，也不是来自中国的创投，而是来自
中东的主权财富基金。2024--2026年，以Abu Dhabi的MGX和沙特的PIF为代表
的主权资本迅速成为全球AI融资中不可忽视的力量。他们的入局不仅是金融行为，
更是国家战略。

\subsection{MGX：Abu Dhabi的AI超级基金}
\label{subsec:inv-mgx}

\begin{corebox}[机构档案：MGX]
\textbf{全称}：MGX（Abu Dhabi人工智能与先进技术委员会下属投资公司） \\
\textbf{成立}：2024年3月 \\
\textbf{背景}：Abu Dhabi主权财富生态的一部分，G42和Mubadala为创始合作伙伴 \\
\textbf{目标AUM}：1000亿美元 \\
\textbf{核心人物}：Sheikh Tahnoon bin Zayed Al Nahyan（阿联酋国家安全顾问） \\
\textbf{投资范围}：AI基础设施、半导体、AI核心技术和应用
\end{corebox}

MGX的崛起速度令人震惊。从2024年3月成立到2025年底，MGX参与了多笔全球
最大的AI交易：

\textbf{AI基础设施合作伙伴关系。}2024年9月，BlackRock、Global
Infrastructure Partners（GIP）、Microsoft和MGX共同宣布了Global AI
Infrastructure Investment Partnership（GAIIP）——一个目标高达1000亿美元
的AI基础设施投资计划。BlackRock计划筹集300亿美元的权益投资，另外700亿
可来自杠杆债务融资。Microsoft是普通合伙人，NVIDIA担任技术顾问。MGX在
这个合作关系中的角色是提供主权资本——这种资本的特征是期限极长
（主权基金通常以十年甚至数十年为周期评估回报）、政治附加条件最少
（相比美国LP的ESG和中国合规要求）。

\textbf{Stargate项目参与。}MGX是Stargate项目（5000亿美元AI基础设施计划）
的四个核心参与者之一，与OpenAI、SoftBank和Oracle并列。

\textbf{数据中心收购。}2025年10月，MGX参与了一个由NVIDIA、Microsoft、
BlackRock和xAI组成的投资财团，以400亿美元收购了Aligned Data Centers
——这是数据中心历史上最大的收购交易之一。

\textbf{Anthropic Series G。}2026年2月，MGX作为联合领投方之一参与了
Anthropic的300亿美元Series G融资。

\textbf{TikTok美国业务。}2025年9月，MGX还参与了Oracle和Silver Lake主导的
TikTok美国业务收购——虽然不直接是AI交易，但反映了MGX在全球科技投资中
的广泛布局。

\begin{whybox}[Abu Dhabi为什么在AI上如此激进？]
Abu Dhabi通过MGX的AI投资可以从三个层面理解：
\begin{enumerate}
  \item \textbf{经济转型}：阿联酋的石油收入终将下降。AI代表了一种全新的
    经济引擎——如果Abu Dhabi能成为全球AI基础设施的关键节点（通过数据中心、
    算力和能源供应），它就能在后石油时代保持经济影响力。
  \item \textbf{地缘优势}：Abu Dhabi位于美国和中国之间的"灰色地带"。
    当中美AI投资互相封锁时，Abu Dhabi可以同时与双方合作——这种"中间人"
    角色在地缘政治紧张时期价值极高。
  \item \textbf{能源-算力协同}：AI需要巨量电力。中东拥有廉价的化石能源
    和大片沙漠（可建设太阳能电站）。将能源优势转化为算力优势是一种
    自然的产业逻辑——数据中心需要的正是Abu Dhabi丰富的能源供应。
\end{enumerate}
Sheikh Tahnoon bin Zayed Al Nahyan——阿联酋国家安全顾问、MGX的幕后推手
——2025年3月在白宫会见了特朗普总统，讨论AI基础设施合作。这次会面
标志着AI投资已正式成为国家层面外交议程的一部分。
\end{whybox}


\subsection{沙特PIF与HUMAIN：石油财富的AI转化}
\label{subsec:inv-pif}

\begin{corebox}[机构档案：沙特公共投资基金（PIF）]
\textbf{全称}：Public Investment Fund（PIF） \\
\textbf{管理资产}：超过1万亿美元 \\
\textbf{领导人}：Yasir Al-Rumayyan（总裁），穆罕默德·本·萨勒曼王储（董事会主席） \\
\textbf{AI战略载体}：HUMAIN（2025年5月成立） \\
\textbf{关键合作}：Google Cloud AI Hub（2024年10月签署）
\end{corebox}

沙特的AI投资雄心甚至超过了Abu Dhabi——考虑到PIF超过1万亿美元的资产
规模，这并不令人意外。

\textbf{HUMAIN的成立。}2025年5月12日，穆罕默德·本·萨勒曼王储亲自宣布
成立HUMAIN——一家PIF全资持有的AI公司。HUMAIN由王储本人担任董事长，
将覆盖整个AI价值链：下一代数据中心、AI基础设施和云能力、先进AI模型
和解决方案。值得注意的是，HUMAIN还将开发"世界上最强大的多模态阿拉伯语
大语言模型之一"——这标志着AI投资不仅关乎财务回报，更关乎语言和文化主权。

\textbf{Google Cloud合作。}2024年10月，PIF与Google Cloud签署了战略合作，
在沙特东部省份达曼附近建设先进AI Hub。预计在未来八年为沙特GDP贡献710亿
美元，并创造数千个就业机会。

\textbf{Anthropic Series G。}Qatar Investment Authority（QIA）作为重要
投资者参与了Anthropic的\$30B Series G融资，沙特相关基金也通过多个渠道
参与了全球AI融资。

2025年10月，Bloomberg报道PIF将在2026--2030年的投资战略中进一步聚焦
AI相关公司——特别是HUMAIN——目标是将部分投资组合公司培育为"全球冠军"。


\subsection{其他主权参与者}
\label{subsec:inv-other-sovereign}

\paragraph{新加坡GIC}
GIC在2026年2月联合领投了Anthropic的\$30B Series G融资——这标志着新加坡
作为AI投资中心的地位进一步巩固。GIC的投资风格以"长期、稳健、全球分散"
著称，其参与Anthropic的领投在一定程度上为这轮融资提供了"非美国、非中东"
的平衡因素。

\paragraph{卡塔尔投资局（QIA）}
QIA是Anthropic多轮融资的重要参与者。2025年，QIA在Anthropic \$13B Series F
中的参与金额据报道超过10亿美元。QIA的AI投资策略与其更广泛的科技投资
组合一致——通过财务投资获取AI技术的接触权。

\paragraph{日本政府养老投资基金（GPIF）}
虽然GPIF本身不直接投资AI创业公司，但作为全球最大的养老金基金（管理资产
超过1.6万亿美元），它通过外部管理人间接参与了AI投资——包括通过配置给
Sequoia、a16z等顶级VC基金的资金。

\begin{keybox}[主权财富基金AI投资对比（2024--2026/3）]
\small
\begin{tabularx}{\textwidth}{l r l X}
\toprule
\rowcolor{figLightGray}
\textbf{基金} & \textbf{估计AI配置} & \textbf{所在国}
  & \textbf{代表性AI交易} \\
\midrule
MGX          & $\sim$\$50B+   & 阿联酋
  & Stargate, GAIIP, Anthropic G, Aligned DC \\
PIF/HUMAIN   & $\sim$\$30B+   & 沙特
  & Google Cloud Hub, HUMAIN全资 \\
GIC          & $\sim$\$15B+   & 新加坡
  & Anthropic G联合领投 \\
QIA          & $\sim$\$10B+   & 卡塔尔
  & Anthropic F/G, 多笔AI基础设施 \\
Temasek      & $\sim$\$5B+    & 新加坡
  & Anthropic G参投 \\
\bottomrule
\end{tabularx}

\medskip
\noindent 注：以上数据为基于公开报道的估算，实际金额可能更高。
主权基金的许多投资通过子基金和SPV执行，不一定完全公开披露。
\end{keybox}


% =============================================================
\section{资本的逻辑：AI投资的结构性分析}
\label{sec:inv-analysis}
% =============================================================

在梳理了主要人物和机构后，让我们退后一步，从结构性角度分析AI投资生态
的深层逻辑。

\subsection{三层投资逻辑的叠加}
\label{subsec:inv-three-layers}

2024--2026年的AI超级轮次之所以能达到史无前例的规模，根本原因在于三层
完全不同的投资逻辑在同一个市场中叠加：

\textbf{第一层：财务投资。}Thrive Capital、Lightspeed、a16z、Sequoia等
传统VC基金遵循经典的风险投资逻辑——投入资本、获取股权、等待退出。
他们面临的核心挑战是：当OpenAI估值8400亿、Anthropic估值3800亿时，
这些公司需要在未来十年创造数千亿美元的收入，投资者才能获得合理回报。

\textbf{第二层：战略投资。}Amazon投500亿给OpenAI，核心回报不是股权
增值，而是AWS上锁定最强AI模型的优先部署权。NVIDIA投300亿给OpenAI，
换来的是最大GPU客户的深度绑定。Microsoft早年对OpenAI的投资已带来
数百亿美元的Azure AI收入增长。对这些战略投资者来说，投资回报的计算
公式完全不同于财务VC。

\textbf{第三层：主权战略投资。}GIC、MGX、QIA等主权基金投资AI的逻辑
超越了商业——它关乎国家在AI时代的战略定位。对阿联酋和沙特来说，
即使投资的财务回报不理想，获得的AI技术接触权、产业关系和人才网络
也具有不可替代的战略价值。

三层逻辑的叠加解释了一个核心悖论：\textbf{用传统财务指标衡量，很多
AI投资看起来"不合理"——因为超过一半的资本根本不是按财务回报的逻辑
配置的}。

\subsection{Lightspeed——Anthropic的忠实伙伴}
\label{subsec:inv-lightspeed}

在众多AI VC中，Lightspeed Venture Partners与Anthropic的关系值得特别
分析——它展示了"先发优势"和"持续加注"如何在AI投资中产生巨大回报。

Lightspeed于2023年初首次接触Anthropic——当时这家公司不到100名员工、
没有公开产品、零收入。仅仅两年后，Anthropic已发布业界领先的模型，
建立了与科技巨头的重要合作关系，并以惊人的速度扩张。

\textbf{融资时间线：}
\begin{itemize}
  \item \textbf{2025年3月}：Lightspeed领投Anthropic \$3.5B Series E
    ——当时Anthropic估值约615亿美元；
  \item \textbf{2025年Q3}：领投Anthropic \$4B Series F
    ——估值升至1830亿美元；
  \item \textbf{2026年2月}：再次投资Anthropic \$30B Series G
    ——估值跃至3800亿美元。
\end{itemize}

从Series E到Series G，Lightspeed的早期投资在不到一年内账面增值约6倍。
Lightspeed的Ravi Mhatre、Guru Chahal和Sebastian Duesterhoeft领导了与
Anthropic的合作关系。Earthian AI的2026年AI VC排名将Lightspeed列为
硅谷AI VC的第一名。

\subsection{Coatue、General Catalyst与其他关键参与者}
\label{subsec:inv-others}

\paragraph{Coatue Management}
Philippe Laffont创立的Coatue是一家以数据驱动和技术分析著称的对冲/风投
混合基金。Coatue联合领投了Anthropic \$30B Series G，同时也是多笔AI
基础设施交易的参与者。

\paragraph{General Catalyst}
Hemant Taneja领导的General Catalyst管理超过430亿美元的资产。2024年10月，
Taneja宣布筹集约80亿美元的新基金。General Catalyst的AI投资策略被称为
"Creation Strategy"——不仅投资AI公司，还通过其运营团队帮助被投公司
在组织结构和go-to-market层面进行AI转型。

2026年2月，Taneja在印度AI Impact Summit上宣布对印度科技和创业生态投入
50亿美元——这是对印度AI市场最大的单一VC承诺之一。他的核心论点是：
"印度有580万IT从业者，是全球最大的IT出口国。那些大规模构建全球IT交付
体系的工程师和运营者，现在最有能力直接面对客户并在应用层驱动创新。"

\paragraph{Accel}
Accel参与了Anthropic的Series G，同时在AI DevTools领域有广泛布局。
作为最早投资Facebook的基金之一，Accel的AI投资更侧重于企业级AI应用
和开发者工具。

\paragraph{IVP与Spark Capital}
IVP（Institutional Venture Partners）和Spark Capital在AI投资中扮演
"稳健跟投者"的角色——它们通常不领投最大的AI轮次，但在多家AI公司的
cap table上出现。Spark Capital早期投资了Twitter和Slack等公司，在AI时代
则布局了多个AI应用公司。


% =============================================================
\section{AI投资人物图谱}
\label{sec:inv-landscape}
% =============================================================

综合全章分析，我们可以将AI投资人物按其功能角色进行系统性分类。

\begin{keybox}[AI投资人物功能分类（2024--2026）]
\small
\begin{tabularx}{\textwidth}{l l X}
\toprule
\rowcolor{figLightGray}
\textbf{角色} & \textbf{代表人物} & \textbf{核心功能} \\
\midrule
超级赌徒     & 孙正义 &
  极端集中押注，以国家级规模配置资本 \\
叙事塑造者   & Marc Andreessen, Martin Casado &
  定义AI投资的认知框架和话语体系 \\
关系网络者   & Josh Kushner, Reid Hoffman &
  通过密切的创始人关系获取优先配额和信息 \\
黑天鹅猎人   & Vinod Khosla, Elad Gil &
  在共识形成前下注，捕获100倍以上回报 \\
系统化玩家   & Pat Grady (Sequoia), Lightspeed &
  用纪律化的流程持续发现和支持AI赢家 \\
算力守门人   & Anjney Midha &
  通过GPU分配权影响创业公司生死 \\
批量孵化者   & Garry Tan (YC), 陆奇 (奇绩创坛) &
  通过加速器模式批量生产AI创业公司 \\
三栖人物     & 李开复, Eric Schmidt &
  同时跨越研究、投资和创业/政策三个领域 \\
实用主义者   & 朱啸虎, 张磊 &
  关注AI的实际商业化而非概念炒作 \\
桥梁建造者   & 沈南鹏 &
  在中美科技脱钩趋势中维持跨境资本通道 \\
主权操盘手   & Sheikh Tahnoon, MBS &
  将石油财富转化为AI时代的国家竞争力 \\
政策-投资混合体 & Eric Schmidt, Peter Thiel &
  将投资决策与国家安全和政策框架绑定 \\
\bottomrule
\end{tabularx}
\end{keybox}


\begin{figure}[htbp]
\centering
\begin{tikzpicture}[
  node distance=0.8cm and 1.2cm,
  every node/.style={font=\small},
  tier/.style={draw, rounded corners=4pt, minimum width=3cm,
    minimum height=0.7cm, align=center, font=\small\bfseries},
  t1/.style={tier, fill=figBlue!15, draw=figBlue},
  t2/.style={tier, fill=figOrange!15, draw=figOrange},
  t3/.style={tier, fill=figGreen!15, draw=figGreen},
  t4/.style={tier, fill=figRed!15, draw=figRed},
  arr/.style={-{Stealth[length=4pt]}, thick, figGray},
]

% Title
\node[font=\normalsize\bfseries] at (5.5, 5.5) {AI资本生态的四层结构};

% Layer 1 — Sovereign
\node[t4, minimum width=12cm] (sov) at (5.5, 4.5)
  {主权资本：MGX, PIF/HUMAIN, GIC, QIA};

% Layer 2 — Strategic
\node[t2, minimum width=12cm] (str) at (5.5, 3.3)
  {科技巨头战略投资：Amazon \$50B, NVIDIA \$30B, Microsoft};

% Layer 3 — Mega VC
\node[t1, minimum width=12cm] (mvc) at (5.5, 2.1)
  {顶级VC：a16z, Sequoia, Thrive, Lightspeed, Founders Fund};

% Layer 4 — Early/Angel
\node[t3, minimum width=12cm] (ang) at (5.5, 0.9)
  {早期/天使：Khosla, Elad Gil, YC, Gross\&Friedman, 奇绩创坛};

% Arrows showing capital flow
\draw[arr] (sov.south) -- (str.north);
\draw[arr] (str.south) -- (mvc.north);
\draw[arr] (mvc.south) -- (ang.north);

% Side label
\node[font=\scriptsize, text=figGray, rotate=90, anchor=south] at (-0.8, 2.7)
  {单笔规模递减 $\leftarrow$};
\node[font=\scriptsize, text=figGray, rotate=90, anchor=north] at (11.8, 2.7)
  {$\rightarrow$ 投资数量递增};

\end{tikzpicture}
\caption{AI资本生态的四层结构：从主权资本到天使投资}
\label{fig:inv-capital-layers}
\end{figure}


% =============================================================
\section{本章总结：资本如何塑造AI的未来？}
\label{sec:inv-conclusion}
% =============================================================

\begin{corebox}[本章核心要点总结]
\begin{enumerate}
  \item \textbf{规模史无前例}：2024--2026年的AI融资规模超越了任何
    历史类比。OpenAI单轮1100亿、Anthropic累计超300亿、Stargate项目
    5000亿——这些数字不是"又一个融资周期"，而是全新的资本形态。

  \item \textbf{少数人的影响力}：大约50个关键人物——本章所介绍的投资者
    和决策者——控制或影响了全球AI投资的大部分资本配置。他们的认知框架、
    关系网络和风险偏好，直接决定了哪些AI公司能获得资源，哪些将被边缘化。

  \item \textbf{三层逻辑叠加}：财务投资、战略投资和主权投资的叠加
    创造了AI融资的天文数字。用单一维度理解AI估值会产生误导——超过一半
    的资本不是按传统财务回报逻辑配置的。

  \item \textbf{地缘政治化}：AI投资已不再是纯粹的商业行为。中美投资
    脱钩、中东主权基金入局、Stargate项目在白宫宣布——资本流向正在被
    地缘政治深刻重塑。

  \item \textbf{算力即权力}：Anjney Midha和a16z的Oxygen项目揭示了一个
    新的权力维度——在GPU稀缺时代，算力分配权可能比资本分配权更重要。
    这改变了传统VC的权力结构。

  \item \textbf{投资人自身的被投化}：Daniel Gross和Nat Friedman被Meta
    收购的案例显示，在AI时代，优秀的投资人——因其关系网和技术洞察力
    ——本身成为了被争抢的资产。

  \item \textbf{中国的独立路径}：沈南鹏、李开复、张磊、朱啸虎、陆奇
    等中国投资者在地缘政治约束下走出了差异化道路——更强调应用层、
    更深的政府合作、更本地化的投资逻辑。

  \item \textbf{泡沫信号}：2026年3月，Benchmark的Bill Gurley公开警告
    "AI泡沫即将破裂"，First Round Capital的Howard Morgan表示
    "AI估值已过热"。资本的疯狂涌入是否可持续，是未来12--18个月
    最关键的悬念之一。
\end{enumerate}
\end{corebox}

在AI时代的宏大叙事中，资本扮演着一个矛盾的角色：它既是AI进步的
必要条件（没有数百亿美元的投入，就不会有GPT-4、Claude、Gemini），
又可能是AI泡沫的催化剂（当FOMO驱动的资本追逐有限的优质标的，估值
必然脱离基本面）。

本章介绍的50位投资人和决策者，在某种意义上承担着一种独特的责任：
他们不仅需要为自己的LP创造回报，还在事实上参与决定AI技术将朝什么
方向发展、以什么速度推进、受什么约束条件限制。当我们讨论AI的未来时，
理解这些"资本推手"的思维方式和利益结构，与理解技术本身同样重要。

下一章（\S\ref{ch:founders}）将聚焦AI创业公司的创始人——
那些实际使用这些资本来构建产品和公司的人。如果说本章介绍的是"油门"，
下一章介绍的就是"方向盘"。
